\documentclass[11.5pt]{article}

\usepackage[T1]{fontenc}
\usepackage[utf8]{inputenc}
\usepackage[spanish]{babel}

\title{Proyecto Intermodular: Opalo E-Learning}
\author{Iván Gimeno Montoya e Izan Gregorio Vizcarra}
\date{17 de mayo de 2026}

\begin{document}
	
	\maketitle
	
	\newpage

	\tableofcontents
	
	\newpage
	\section{Objetivos de la práctica}
	
	La realización de esta memoria tiene como objetivo analizar de forma integral el Proyecto Intermodular para entender su contexto, diseño y solución técnica, planificación, implementación y desarrollo técnico, verificación de su funcionamiento, análisis de sostenibilidad y valoración personal del mismo.
	
	\section{Relación de herramientas (soft/hard) utilizadas en el desarrollo de la práctica.}
	
	\begin{itemize}
		\item Ordenador
		\item Con acceso a:
		\subitem Las prácticas anteriores
		\subitem Internet para la investigación
	\end{itemize}
	
	\section{Realización de la práctica}
	
	
	\subsection{Introducción y Contextualización}
	
	\subsubsection{Título del proyecto}
	Opalo E-Learning.
	\subsubsection{Justificación}
	Opalo E-Learning es relevante como herramienta en el contexto actual porque tiene como único objetivo la formación, desde la introducción hasta la comprensión más profunda del ámbito de conocimiento a tratar, en temas de interés para nuestros usuarios. Se diferencia de otras plataformas de aprendizaje porque aquí: no tratamos de vender humo, toda la formación es enteramente gratuita, y nos basamos en fuentes fiables y revisables.
	
	\subsubsection{Objetivos generales}
	El objetivo general del proyecto es diseñar, desarrollar y desplegar una página web corporativa, funcional, con diseño responsive o adaptable, y accesible para todo el mundo para una empresa que ofrece servicios de formación, integrando las competencias de módulos de SMR2. 
	
	Objetivos específicos que se reconocen:
	\begin{enumerate}
		\item Diseñar una página web de forma responsiva y accesible para todos.
		
		\item Implementarla mediante herramientas como Apache, XAMPP y Acrylic DNS Proxy.
		
		\item Crear las presentaciones de H5P para la creación de los cursos.
		
		\item Documentar todo el proceso desde el principio hasta el final.
		
		\item Optimizar el código para que sea lo más ligero posible y compatible con todo tipo de dispositivos, incluso siendo éstos de bajo rendimiento.
		
		\item Por último, explicar a un jurado de forma oral nuestra empresa y página web.
	\end{enumerate}
	
	\subsubsection{Identificación de necesidades}
	
	Hemos detectado que la realidad de muchas webs del sector tecnológico (sector formativo) no cumplen con tres aspectos clave:
	
	1. Diseño y desarrollo web:
	
	Numerosas compañías eligen emplear plantillas estándar de WordPress o Wix, mientras que algunas otras ni siquiera poseen una página web profesional. Aquí proponemos un sitio web bien diseñado y agradable para el visitante.
	
	2. Educar sobre los temas que se tratan:
	
	Es frecuente que las empresas del mismo sector no proporcionen a sus posibles clientes toda la información necesaria para comprender la relevancia de sus servicios (aunque inicialmente se planificó el proyecto como una empresa del sector tecnológico, actualmente, estamos enfocados en el sector de la formación en Internet).
	Esto reduce la confianza y el interés del posible comprador, pues no comprende a qué se dedica la empresa. Demostrando compromiso con la transparencia, nuestro equipo está a su disposición para proporcionarle toda la información que tenemos. Por esta razón, disponemos de una sección dedicada a nosotros y otra de Contacto.
	
	3. Rendimiento:
	
	Una buena parte de las páginas web no se ajustan a los criterios de una carga veloz de todo el contenido.
	
	
	
	
	
	
	
	
	\section{Conclusiones}
	
	En esta sección se resumen las conclusiones finales del trabajo.
	
\end{document}